\documentclass[11pt]{article}

% ---------- packages ----------
\usepackage[margin=1in]{geometry}
\usepackage{amsmath, amssymb, amsthm}
\usepackage{booktabs}      % nice tables
\usepackage{enumitem}      % control over (a), (b), ...
\usepackage{siunitx}       % \num{} for consistent sig figs
\usepackage[colorlinks=true, urlcolor=blue]{hyperref}

% ---------- shortcuts ----------
\newcommand{\E}{\mathbb{E}}
\newcommand{\Prob}{\mathbb{P}}
\newcommand{\Var}{\operatorname{Var}}
\newcommand{\Cov}{\operatorname{Cov}}
\newcommand{\R}{\mathbb{R}}
\newcommand{\indep}{\perp\!\!\!\perp}
\newcommand{\dd}{\,\mathrm{d}}

% consistent rounding: \num[round-mode=figures,round-precision=4]{0.1234567}
\sisetup{round-mode=figures, round-precision=4}

\title{Markov Processes: Homework 1}
\author{Linden Kurtz}
\date{Due September 2, 2026}

\begin{document}
\maketitle

\noindent
GitHub account: \texttt{lindenkurtz} \\
Repository: \url{https://github.com/lindenkurtz/Markov2026} \\
Handwritten notes can also be found in this repo.

\section*{Problem 1: Joint distributions}

\begin{enumerate}[label=(\alph*)]
  \item
  Due to the nature of probability distributions, $\int_0^2\!\!\int_0^1 K(x^2+y)\dd x \dd y = 1$.
  \begin{align}
    \int_0^1\!\!\int_0^2 x^2\dd y \dd x + \int_0^1\!\!\int_0^2 y\dd y \dd x &= \tfrac{1}{K}, \\
    \int_0^1 2x^2 \dd x + \int_0^1 2 \dd x &= \tfrac{1}{K}, \\
    \tfrac{2}{3} + 2 &= \tfrac{1}{K}
  \end{align}
  Therefore $K = \tfrac{3}{8}$.

  \item
  The marginal of $X$ is $f_X(x) = \int_0^2 f(x,y)\dd y$ for $0 \le x \le 1$, and zero otherwise.
  \[
    f_X(x) = \int_0^2 \tfrac{3}{8} (x^2 + y)\dd y = \tfrac{3}{4}x^2+\tfrac{3}{4}
  \]
  As a check, $\int_0^1 f_X(x)\dd x = \int_0^1 \tfrac{3}{4}x^2+\tfrac{3}{4}\dd x = 1$. \\ \\
  Similarly, the marginal of $Y$ is $f_Y(y) = \int_0^1 f(x,y)\dd x$ for $0 \le y \le 2$, and zero otherwise.
  \[
    f_Y(y) = \int_0^1 \tfrac{3}{8} (x^2 + y)\dd x = \tfrac{1}{8}+\tfrac{3}{8}y
  \]
  As a check, $\int_0^2 f_Y(y)\dd y = \int_0^2 \tfrac{1}{8}+\tfrac{3}{8}y\dd y = 1$.

  \item
  $X$ and $Y$ are independent if and only if $f(x,y) = f_X(x)f_Y(y)$ for all $(x,y)$.
  Comparing the two expressions,
  \begin{align}
    f_X(x)f_Y(y) &= (\tfrac{3}{4}x^2 + \tfrac{3}{4})(\tfrac{1}{8}+\tfrac{3}{8}y), \\
    &= \tfrac{3}{32}x^2 + \tfrac{9}{32}x^2 y + \tfrac{3}{32} + \tfrac{9}{32}y, \\
    &\not= f(x,y) = \tfrac{3}{8}(x^2 + y)
  \end{align}
  So $X$ and $Y$ are not independent.
\end{enumerate}

\section*{Problem 2: Moment generating function 1}

\begin{enumerate}[label=(\alph*)]
  \item Given $X \sim \text{Geometric}(q)$ counting misses, so $\Prob(X = x) = (1-q)^x q$. Therefore, the MGF is given by:
  \begin{align}
    M_X(s) &= \E[e^{Xs}], \\
    &= \sum_{x=0}^{\infty} e^{xs}(1-q)^x q, \\
    &= q \sum_{x=0}^{\infty} (e^s (1-q))^x
  \end{align}
  Assuming that $e^s(1-q) < 1$, or $e^s < \tfrac{1}{1-q}$,
  \[
    M_X(s) = \frac{q}{1-e^s(1-q)}
  \]
  And the MGF is finite for $s < \ln(\tfrac{1}{1-q})$.

  \item Differentiating at $s=0$:
  \begin{align}
    M_X'(s) &= \frac{q(1-q)e^s}{(1-e^s(1-q))^2}, \\
    M_X'(0) &= \frac{1-q}{q}, \\
    M_X''(s) &= \frac{(1-e^s(1-q))^2(q(1-q)e^s) + 2(q(1-q)e^s)(1-e^s(1-q))(1-q)e^s}{(1-e^s(1-q))^4}, \\
    M_X''(0) &= \frac{(1-q)(q+2(1-q))}{q^2}, \\
    \E[X] &= M_X'(0) = \frac{1-q}{q}, \\
    \Var(X) &= M_X''(0) - (M_X'(0))^2 = \frac{1-q}{q^2}
  \end{align}

  \item Fix an integer $k \ge 0$. Writing the conditional MGF of $X-k$ as an expectation and using
  $\E[g(X)] = \sum_x g(x)\Prob(X=x)$,
  \begin{align}
    M_{X-k \mid X \ge k}(s) &= \E\!\left[e^{s(X-k)} \mid X \ge k\right], \\
    &= \sum_{n=k}^{\infty} e^{s(n-k)}\,\Prob(X = n \mid X \ge k), \\
    &= \sum_{n=k}^{\infty} e^{s(n-k)}\,\frac{\Prob(X=n \cap X \ge k)}{\Prob(X \ge k)}
  \end{align}
  Every $n$ in this summation already satisfies $n \ge k$, so the event $\{X = n \cap X \ge k\}$
  is just $\{X = n\}$, and $\Prob(X \ge k) = (1-q)^k$. Therefore
  \begin{align}
    M_{X-k \mid X \ge k}(s) &= \sum_{n=k}^{\infty} e^{s(n-k)}\,\frac{(1-q)^n q}{(1-q)^k}, \\
    &= \frac{q\,e^{-sk}}{(1-q)^k}\sum_{n=k}^{\infty} \left(e^s(1-q)\right)^n, \\
    &= \frac{q\,e^{-sk}}{(1-q)^k}\cdot\frac{\left(e^s(1-q)\right)^k}{1-e^s(1-q)}, \\
    &= \frac{q}{1-e^s(1-q)} \;=\; M_X(s)
  \end{align}
  Since the MGF determines the distribution, this means that the number of further misses for a
  player who has already missed $k$ shots has the same distribution as if they had shot no shots at
  all. 

  \item Let $T_r = \sum_{i=1}^{r} X_i$, where the $X_i$ are i.i.d.\ random variables with the same
  distribution as $X$. Since the $X_i$ are independent, the MGF of the sum is the product of the MGFs:
  \begin{align}
    M_{T_r}(s) &= \left(M_X(s)\right)^r = \left(\frac{q}{1-e^s(1-q)}\right)^{r}
  \end{align}
  This is the MGF of a negative binomial distribution counting the number of misses before the $r$th
  make, with success probability $q$ and $r$ successes. Differentiating,
  \begin{align}
    M_{T_r}'(s) &= r\left(q(1-e^s(1-q))^{-1}\right)^{r-1}\cdot\frac{q(1-q)e^s}{(1-e^s(1-q))^2}, \\
    M_{T_r}'(0) &= r(1)^{r-1}\cdot\frac{1-q}{q}, \\
    \E[T_r] &= \frac{r(1-q)}{q}
  \end{align}
  and for the second moment,
  \begin{align}
    M_{T_r}''(0) &= r\cdot\frac{(1-q)(q+2(1-q))}{q^2}
      + \frac{1-q}{q}\cdot\frac{r(r-1)(1-q)}{q}, \\
    &= \frac{r(1-q)(q+2(1-q)) + r(r-1)(1-q)^2}{q^2}
  \end{align}
  so that
  \begin{align}
    \Var(T_r) &= \frac{r(1-q)(q+2(1-q)) + r(r-1)(1-q)^2}{q^2} - \frac{r^2(1-q)^2}{q^2}, \\
    &= \frac{r(1-q)\big(-q+2+(r-1)(1-q)-r(1-q)\big)}{q^2}, \\
    &= \frac{r(1-q)\big(-q+2+r-rq-1+q-r+rq\big)}{q^2}, \\
    &= \frac{r(1-q)}{q^2}
  \end{align}
  Both the mean and the variance are exactly $r$ times the corresponding quantities for a single $X$,
  which is what a sum of $r$ independent copies should give.
\end{enumerate}

\section*{Problem 3: Mixtures and conditioning}

Let $\theta \sim \text{Uniform}[0,1]$ and $X \mid \theta \sim \text{Exponential}(\theta)$.
The joint density follows from $f_{X\mid\theta}(x\mid\theta) = f_{X,\theta}(x,\theta)/f_\theta(\theta)$:
\begin{align}
  f_{X,\theta}(x,\theta) &= f_{X\mid\theta}(x\mid\theta)\,f_\theta(\theta), \\
  &= \theta e^{-\theta x}\cdot 1, \qquad 0 < \theta \le 1,\ x > 0
\end{align}
So, every $0 < \theta \le 1$ contributes to any given $x$. \\

Marginalizing over $\theta$ and integrating by parts with $u = \theta$ and $v' = e^{-\theta x}$
(so $u' = 1$ and $v = -\tfrac{1}{x}e^{-\theta x}$),
\begin{align}
  f_X(x) &= \int_0^1 \theta e^{-\theta x}\dd \theta, \\
  &= \left[-\frac{\theta}{x}e^{-\theta x}\right]_0^1 + \int_0^1 \frac{1}{x}e^{-\theta x}\dd\theta, \\
  &= \left(-\frac{\theta}{x}e^{-\theta x} - \frac{1}{x^2}e^{-\theta x}\right)\Bigg|_0^1, \\
  &= \left(-\frac{1}{x}e^{-x} - \frac{1}{x^2}e^{-x}\right) - \left(0 - \frac{1}{x^2}\right), \\
  &= \frac{1}{x^2} - e^{-x}\left(\frac{1}{x} + \frac{1}{x^2}\right), \qquad x > 0
\end{align}

For the expected value:
\begin{align}
  \E[X] = \E\big[\E[X \mid \theta]\big] = \E\!\left[\frac{1}{\theta}\right]
  = \int_0^1 \frac{1}{\theta}\dd\theta = \infty
\end{align}
so $\E[X]$ is not finite. This aligns with the result from taking the expected value of $f_X(x)$ directly, since an integrand of $\tfrac{1}{x}$ diverges.

\section*{Problem 4: Moment generating function 2}

\begin{enumerate}[label=(\alph*)]
  \item Let $X \sim \text{Exponential}(\lambda)$ with $f(x) = \lambda e^{-\lambda x}$ for $x \ge 0$.
  \begin{align}
    M_X(s) &= \E[e^{sX}], \\
    &= \int_0^\infty e^{sx}\lambda e^{-\lambda x}\dd x, \\
    &= \lambda \int_0^\infty e^{x(s-\lambda)}\dd x
  \end{align}
  Assuming $s < \lambda$, so that the exponent is negative and the integral converges,
  \begin{align}
    M_X(s) &= \lambda\left(\frac{1}{s-\lambda}e^{x(s-\lambda)}\right)\Bigg|_0^\infty
    = \lambda\left(0 - \frac{1}{s-\lambda}\right) = \frac{-\lambda}{s-\lambda}
  \end{align}
  so the MGF is finite for $s < \lambda$. Differentiating,
  \begin{align}
    M_X'(s) &= \frac{\lambda}{(s-\lambda)^2}, & M_X'(0) &= \frac{\lambda}{\lambda^2} = \frac{1}{\lambda} = \E[X], \\
    M_X''(s) &= \frac{-2\lambda}{(s-\lambda)^3}, & M_X''(0) &= \frac{2}{\lambda^2},
  \end{align}
  and therefore $\Var(X) = \dfrac{2}{\lambda^2} - \dfrac{1}{\lambda^2} = \dfrac{1}{\lambda^2}$.

  \item Let $S_n = X_1 + \dots + X_n$ with $X_i$ i.i.d.\ $\text{Exponential}(\lambda)$. By independence,
  \begin{align}
    M_{S_n}(s) &= \left(\frac{-\lambda}{s-\lambda}\right)^{n}, \\
    M_{S_n}'(s) &= n\left(\frac{-\lambda}{s-\lambda}\right)^{n-1}\cdot\frac{\lambda}{(s-\lambda)^2}, \\
    M_{S_n}'(0) &= n(1)^{n-1}\cdot\frac{1}{\lambda} = \frac{n}{\lambda} = \E[S_n], \\
    M_{S_n}''(s) &= n(n-1)\left(\frac{-\lambda}{s-\lambda}\right)^{n-2}\cdot\frac{\lambda}{(s-\lambda)^2}\cdot\frac{\lambda}{(s-\lambda)^2}
      + n\left(\frac{-\lambda}{s-\lambda}\right)^{n-1}\cdot\frac{-2\lambda}{(s-\lambda)^3}, \\
    M_{S_n}''(0) &= n(n-1)\cdot\frac{1}{\lambda^2} + n\cdot\frac{2}{\lambda^2}
      = \frac{n(n-1)+2n}{\lambda^2} = \frac{n(n+1)}{\lambda^2}, \\
    \Var(S_n) &= \frac{n(n+1)}{\lambda^2} - \frac{n^2}{\lambda^2} = \frac{n}{\lambda^2}
  \end{align}
  This MGF is that of a gamma distribution with shape $n$ and rate $\lambda$, so the total waiting
  time for $n$ exponential events is $\text{Gamma}(n, \lambda)$.
\end{enumerate}

\section*{Problem 5: Luck, Skill, and Upsets}

\begin{enumerate}[label=(\alph*)]
  \item With $\sigma(\beta s) = \dfrac{1}{1+e^{-\beta s}}$ and $\beta > 0$, the match is a coin flip with
  probability $\alpha$ and decided by skill otherwise, so
  \begin{align}
    f(s) &= \alpha\left(\tfrac{1}{2}\right) + (1-\alpha)\sigma(\beta s)
    = \frac{\alpha}{2} + (1-\alpha)\left(\frac{1}{1+e^{-\beta s}}\right)
  \end{align}
  Replacing $s$ with $-s$ and multiplying through by $e^{-\beta s}/e^{-\beta s}$,
  \begin{align}
    f(-s) &= \frac{\alpha}{2} + (1-\alpha)\left(\frac{1}{1+e^{\beta s}}\right)\cdot\frac{e^{-\beta s}}{e^{-\beta s}}, \\
    &= \frac{\alpha}{2} + (1-\alpha)\left(\frac{e^{-\beta s}}{1+e^{-\beta s}}\right), \\
    &= \frac{\alpha}{2} + (1-\alpha)\left(\frac{1+e^{-\beta s}}{1+e^{-\beta s}} - \frac{1}{1+e^{-\beta s}}\right), \\
    &= \frac{\alpha}{2} + (1-\alpha) - (1-\alpha)\left(\frac{1}{1+e^{-\beta s}}\right), \\
    &= 1 - \frac{\alpha}{2} - (1-\alpha)\left(\frac{1}{1+e^{-\beta s}}\right) = 1 - f(s)
  \end{align}
  For the limits,
  \begin{align}
    \lim_{s\to\infty} f(s) &= \frac{\alpha}{2} + (1-\alpha) = 1 - \frac{\alpha}{2}, &
    \lim_{s\to-\infty} f(s) &= \frac{\alpha}{2}
  \end{align}
  As player $i$ gets infinitely better than player $j$, the chance he wins approaches $1$ minus the
  chance the game is determined by luck and lost on the coin flip. As player $i$ gets infinitely
  worse, the only way he wins is on the coin flip, so his chance of winning is the chance the game is
  determined by luck times one half.

  \item With $\alpha = 0.2$, $\beta = 1$, and $s = 2$,
  \begin{align}
    \Prob(i \text{ wins}) &= f(2) = \frac{0.2}{2} + (1-0.2)\left(\frac{1}{1+e^{-2}}\right)
    = \frac{1}{10} + \frac{4}{5}\left(\frac{1}{1+e^{-2}}\right) \approx \num{0.8046}
  \end{align}
  For a \$1 wager at even odds the bettor receives \$2 on a win and nothing on a loss, so with
  $W$ the amount returned,
  \begin{align}
    \E[W] &= 2\,\Prob(i \text{ wins}) + 0\cdot\left(1 - \Prob(i \text{ wins})\right)
    = 2(\num{0.8046}) \approx \num{1.609}
  \end{align}
  That figure includes the \$1 that was staked, so the expected net return is
  $\num{1.609} - 1 \approx \$\num{0.6092}$ per game. Finally, letting $C$ be the event that the coin
  decided the match, we have $\Prob(C) = 0.2$ and
  $\Prob(j \text{ wins}) = 1 - \num{0.8046} \approx \num{0.1954}$, so Bayes' theorem gives
  \begin{align}
    \Prob(C \mid j \text{ wins}) &= \frac{\Prob(j \text{ wins} \mid C)\,\Prob(C)}{\Prob(j \text{ wins})}
    = \frac{(0.5)(0.2)}{\num{0.1954}} \approx \num{0.5119}
  \end{align}

  \item Let $E$ be the event that the weaker player won and $F$ the event that the match was decided
  by the coin, with $s > 0$. Then
  \begin{align}
    \Prob(E) &= 1 - f(s) = \frac{\alpha}{2} + (1-\alpha)\frac{1}{1+e^{\beta s}}, \\
    \Prob(F) &= \alpha, \\
    \Prob(E \mid F) &= \frac{1}{2}
  \end{align}
  The two events are independent exactly when $\Prob(E \mid F) = \Prob(E)$, so $\alpha$ and $s$ must
  be chosen so that $\Prob(E) = \tfrac{1}{2}$:
  \begin{align}
    \frac{1}{2} &= \frac{\alpha}{2} + (1-\alpha)\frac{1}{1+e^{\beta s}}, \\
    0 &= \frac{\alpha-1}{2} + (1-\alpha)\frac{1}{1+e^{\beta s}}, \\
    0 &= (1+e^{\beta s})(\alpha-1) + 2(1-\alpha), \\
    &= \alpha + \alpha e^{\beta s} - 1 - e^{\beta s} + 2 - 2\alpha, \\
    0 &= -\alpha + \alpha e^{\beta s} - e^{\beta s} + 1, \\
    0 &= (\alpha - 1)\left(e^{\beta s} - 1\right)
  \end{align}
  The product vanishes when $\alpha = 1$ or when $s = 0$, but since the problem fixes $s > 0$, the
  second root is unavailable and we only take $\alpha = 1$. So $E$ and $F$ are independent when $(\alpha, s) = (1, s)$ with $s > 0$, and are dependent for every other pair.

  \item Expanding $f$ about $s = 0$ with the Maclaurin series
  $f(s) = \sum_{i=0}^{\infty} \frac{f^{(i)}(0)}{i!}s^i$:
  \begin{align}
    f(0) &= \frac{1}{2}, \\
    f'(s) &= -(1-\alpha)\frac{1}{(1+e^{-\beta s})^2}\cdot\left(-\beta e^{-\beta s}\right)
      = (1-\alpha)\frac{\beta e^{-\beta s}}{(1+e^{-\beta s})^2}, \\
    f'(0) &= (1-\alpha)\frac{\beta}{4}
  \end{align}
  and for the second derivative,
  \begin{align}
    f''(s) &= (1-\alpha)\frac{(1+e^{-\beta s})^2\left(-\beta^2 e^{-\beta s}\right)
      - \beta e^{-\beta s}\cdot 2(1+e^{-\beta s})\left(-\beta e^{-\beta s}\right)}{(1+e^{-\beta s})^4}, \\
    &= (1-\alpha)\frac{-\beta^2 e^{-\beta s}(1+e^{-\beta s}) + 2\beta^2 e^{-2\beta s}}{(1+e^{-\beta s})^3}, \\
    &= (1-\alpha)\frac{\beta^2 e^{-\beta s}\left(e^{-\beta s} - 1\right)}{(1+e^{-\beta s})^3}, \\
    f''(0) &= 0
  \end{align}
  The quadratic term vanishes, so
  \begin{align}
    f(s) &= \frac{1}{2} + (1-\alpha)\frac{\beta}{4}\,s + O(s^3)
  \end{align}
  A luck factor $\alpha$ doesn't play a large role across a
  large number of evenly matched games becuase as $s \to 0$, the terms containing $\alpha$ also go to zero, and the only major term remaining is 
  constant $\tfrac{1}{2}$, which doesn't depend on $\alpha$.

  As $s$ increases or decreases, the luck factor $\alpha$ becomes more influential, eventually being
  the only parameter that matters as $s \to \pm\infty$. So, more lopsided matches break this tie.
\end{enumerate}

\section*{Problem 6: Simulation Assignment 1: Coupled Inequalities}

\begin{enumerate}[label=(\alph*)]
  \item
  Let $X, Y, Z$ be independent $\text{Uniform}[0,1]$. The two conditions are $X^2 + Y^2 < Z$ and
  $Z^2 > XY$. Since both are lower bounds
  on $z$ with no upper bound other than $z \le 1$, so the region is the set of points lying above
  whichever of the two surfaces is larger.

  Working in cylindrical coordinates with $x = r\cos\theta$, $y = r\sin\theta$ and
  $\theta \in [0, \tfrac{\pi}{2}]$ (since $x, y \ge 0$), the two bounds become
  \begin{align}
    z > r^2, \qquad z > \sqrt{r^2\cos\theta\sin\theta} = r\,\frac{\sqrt{\sin(2\theta)}}{\sqrt{2}}
  \end{align}
  They exchange dominance where $r^2 = r\sqrt{\cos\theta\sin\theta}$, that is at
  \begin{align}
    r^*(\theta) = \sqrt{\cos\theta\sin\theta} = \frac{\sqrt{\sin(2\theta)}}{\sqrt{2}}
  \end{align}
  For $r > r^*$ the parabolic bound $r^2$ is the larger one, and for $r < r^*$ the $\sqrt{xy}$ bound
  is. Splitting the radial integral at $r^*$ and including the Jacobian $r$,
  \begin{align}
    P &= \int_0^{\pi/2}\!\!\int_{\frac{\sqrt{\sin 2\theta}}{\sqrt{2}}}^{1}\!\!\int_{r^2}^{1} r \dd z \dd r \dd\theta
      + \int_0^{\pi/2}\!\!\int_{0}^{\frac{\sqrt{\sin 2\theta}}{\sqrt{2}}}\!\!
        \int_{r\frac{\sqrt{\sin 2\theta}}{\sqrt{2}}}^{1} r \dd z \dd r \dd\theta, \\
    &= \int_0^{\pi/2}\!\!\int_{\frac{\sqrt{\sin 2\theta}}{\sqrt{2}}}^{1} \left(r - r^3\right) \dd r \dd\theta
      + \int_0^{\pi/2}\!\!\int_{0}^{\frac{\sqrt{\sin 2\theta}}{\sqrt{2}}}
        \left(r - r^2\frac{\sqrt{\sin 2\theta}}{\sqrt{2}}\right) \dd r \dd\theta
  \end{align}
  Carrying out the radial integrals,
  \begin{align}
    P &= \int_0^{\pi/2}\left[\left(\frac{1}{2}-\frac{1}{4}\right)
      - \left(\frac{\sin(2\theta)}{4} - \frac{\sin^2(2\theta)}{16}\right)\right]\dd\theta
      + \int_0^{\pi/2}\left[\frac{\sin(2\theta)}{4} - \frac{\sin^2(2\theta)}{12}\right]\dd\theta
  \end{align}
  Using $\sin^2(2\theta) = \tfrac{1}{2}\left(1-\cos(4\theta)\right)$ and clearing denominators,
  \begin{align}
    P &= \int_0^{\pi/2}\frac{\frac{1-\cos(4\theta)}{2} - 4\sin(2\theta) + 4}{16}\dd\theta
      + \int_0^{\pi/2}\frac{3\sin(2\theta) - \frac{1-\cos(4\theta)}{2}}{12}\dd\theta, \\
    &= \int_0^{\pi/2}\frac{-\cos(4\theta) - 8\sin(2\theta) + 9}{32}\dd\theta
      + \int_0^{\pi/2}\frac{6\sin(2\theta) - 1 + \cos(4\theta)}{24}\dd\theta
  \end{align}
  Evaluating the integrals results in:
  \begin{align}
    P &= \left(\frac{-\sin(4\theta)}{128} + \frac{\cos(2\theta)}{8} + \frac{9\theta}{32}\right)\Bigg|_0^{\pi/2}
      + \left(\frac{-\cos(2\theta)}{8} - \frac{\theta}{24} + \frac{\sin(4\theta)}{96}\right)\Bigg|_0^{\pi/2}, \\
    &= \left(0 + \frac{-1}{8} + \frac{9\pi}{64}\right) - \frac{1}{8}
      + \left(\frac{1}{8} - \frac{\pi}{48}\right) + \frac{1}{8}, \\
    &= \frac{9\pi}{64} - \frac{\pi}{48} = \frac{23\pi}{192} \approx \num{0.3763}
  \end{align}
  So the probability that both inequalities hold is $\tfrac{23\pi}{192} \approx \num{0.3763}$.
\end{enumerate}

\end{document}